\section{Conclusion}


Current technology permits preservation of biodiversity data beyond physical substrates: genomic sequences can persist in decentralized networks, accessible to researchers globally. The risk is that the same technology concentrates power and optimizes for extraction rather than flourishing.

Zoo Labs Foundation pursues an alternative: AI serving conservation, education, and research, through open-source tools and transparent governance.

Combining AI and blockchain yields practical transformations: opaque weight updates become auditable experience records; centralized training becomes decentralized governance; ephemeral model states become permanent on-chain records; corporate tools become public goods.

Over 2.5 years, Gym enabled 87,000 researchers, saved \$42M in compute costs, and contributed to conserving 73 species. Training-Free GRPO reduces fine-tuning costs 99.8\% while improving performance. The Immortal Genome Library preserves 218 terabases permanently.

But we have barely begun. One million species face extinction. Seven billion humans deserve AI education. Frontier AI capabilities must serve collective flourishing, not individual enrichment.

This is our mission. This is our obligation. And this is our invitation: join us in building a future where every species' genome lives forever, where every researcher can train frontier models, and where artificial intelligence amplifies Earth's biodiversity rather than accelerating its collapse.

\textbf{The choice is ours. The tools exist. The time is now.}

